%====================================================================================
%  DOCUMENT PRINCIPAL - TEMPLATE RAPPORT & MÉMOIRE UNIVERSITAIRE
%  Compilation : latexmk -lualatex -outdir=build main.tex
%====================================================================================
\documentclass[
  % draft,             % Activer pour masquer les images et accélérer la compilation
  french,              % Langue principale
  english,             % Langue secondaire
  font=calibri,        % Options de police : calibri, comicsans, standard
  todoNote             % Activer l'affichage des notes et révisions
]{ClasseGuillaume}

% --- MÉTADONNÉES DU DOCUMENT ---
\title{Modèle Universitaire pour Mémoires, Thèses et Rapports d'Ingénierie}
\author{Guillaume \textsc{Winant}}
\date{\today}

% --- BASE BIBLIOGRAPHIQUE ---
\addbibresource{bib/references.bib}

%====================================================================================
%  DÉBUT DU DOCUMENT
%====================================================================================
\begin{document}

%------------------------------------------------------------------------------------
%  PRÉ-CONTENU (FRONTMATTER : Numérotation romaine, sans logo)
%------------------------------------------------------------------------------------
\frontmatter
\masquerLogo % Le logo ne s'affiche pas sur la page de garde ni les préliminaires

% 1. Page de garde principale (en géométrie large)
\geometrieLarge
%------------------------------------------------------------------------------------
%  PAGE DE GARDE PRINCIPALE
%------------------------------------------------------------------------------------
\begin{titlepage}
  \begin{singlespace}
  \centering
  \vspace*{0.04\textheight}

  {\scshape\LARGE \textbf{UCLouvain}\par}
  \vspace{0.4cm}
  {\textsc{\Large Faculté des Sciences}\par}
  \vspace{1.5cm}

  {\textsc{\Large Mémoire de Fin d'Études / Bachelor Thesis}\par}
  \vspace{0.8cm}

  \rule{0.9\linewidth}{0.8pt} \\[0.4cm]
  {\huge\bfseries \thetitle \par} \vspace{0.3cm}
  \rule{0.9\linewidth}{0.8pt} \\[1.5cm]

  \begin{minipage}[t]{0.48\textwidth}
    \begin{flushleft}\large
      \emph{Auteur :}\\
      \textbf{\theauthor}\\
      \footnotesize\texttt{guillaume.winant@student.uclouvain.be}
    \end{flushleft}
  \end{minipage}
  \hfill
  \begin{minipage}[t]{0.48\textwidth}
    \begin{flushright}\large
      \emph{Superviseur :}\\
      \textbf{Prof. Dr. Nom Superviseur}\\
      \footnotesize\texttt{superviseur@uclouvain.be}
    \end{flushright}
  \end{minipage}

  \vfill

  {\large \textit{Mémoire présenté en vue de l'obtention du diplôme de Bachelier}\par}
  \vspace{0.3cm}
  {\textit{Département des Sciences Médicales et Biomédicales}\par}

  \vfill
  \vfill

  {\large \thedate\par}
  \end{singlespace}
\end{titlepage}


% 2. Bascule en géométrie standard ou Word pour les préliminaires
\geometrieWord

% 3. Suivi des révisions (avec cible de retour cliquable pour \todoH)
\begin{singlespace}
\listoftodos[Liste des révisions et tâches en attente]
\label{table_todos}
\end{singlespace}

% 4. Résumé bilingue
%------------------------------------------------------------------------------------
%  RÉSUMÉ & ABSTRACT
%------------------------------------------------------------------------------------
\addchap{Résumé}

Ce modèle LaTeX a été conçu pour structurer des mémoires, thèses et rapports universitaires de haute qualité technique et typographique.
Il propose une mise en page soignée, un découpage modulaire des préambules et des contenus, ainsi qu'une prise en charge des polices modernes OpenType et TrueType via LuaLaTeX.
L'utilisateur peut rapidement personnaliser les métadonnées, la bibliographie, les styles d'en-têtes et les environnements personnalisés (tableaux, figures et illustrations).

\bigskip

\selectlanguage{english}
\addchap{Abstract}

This LaTeX template is designed to structure high-quality academic theses, dissertations, and technical reports.
It provides an elegant modular layout, clean typographical presets, and support for modern OpenType and TrueType fonts using LuaLaTeX.
Users can easily configure metadata, references, custom headers, and floating elements (tables, figures, and custom illustration environments).

\selectlanguage{french}



% 5. Table des matières
\setcounter{tocdepth}{2}
\tableofcontents

%------------------------------------------------------------------------------------
%  CORPS DU TEXTE (MAINMATTER : Numérotation arabe, activation du logo)
%------------------------------------------------------------------------------------
\mainmatter
\afficherLogo % Réactivation du logo pour le corps du document

% --- CHAPITRE I ---
\chapter{Cadre Théorique et Fondements}
\label{chap:fondements}

%------------------------------------------------------------------------------------
%  SECTION 1 : PRÉSENTATION ET TYPOGRAPHIE
%------------------------------------------------------------------------------------
\section{Présentation et Typographie}
\label{sec:presentation}

Ce premier chapitre a pour objectif de démontrer les capacités graphiques et typographiques de ce modèle.
Chaque phrase est rédigée sur une ligne distincte afin de faciliter les revues de code et le versionnement sous Git.

\subsection{Familles de polices disponibles}
\label{subsec:polices}

Le modèle intègre plusieurs polices configurables :
\begin{itemize}
  \item \textbf{Police par défaut :} Calibri (sobre, fluide et parfaitement adaptée aux rapports académiques).
  \item \textbf{Police de titre / accentuation :} {\AsapFont Asap Condensed pour les titres principaux}.
  \item \textbf{Police spécifique :} {\ComicSansMSFont Comic Sans MS} (via \texttt{\textbackslash ComicSansMSFont}).
\end{itemize}

\subsection{Unités et formules scientifiques}
\label{subsec:sciences}

Le package \texttt{siunitx} permet d'écrire des grandeurs physiques sans ambiguïté :
Une vitesse de \qty{25}{\meter\per\second}, une température de \qty{37.2}{\degreeCelsius}, ou encore une masse de \qtyrange{50}{80}{\kilo\gram}.

Les équations mathématiques sont formatées selon les recommandations modernes :
\begin{equation}
  \label{eq:frequence_cardiaque}
  f_{\text{max}} = 220 - \text{âge} \pm \epsilon,
\end{equation}
où $\epsilon$ représente la marge d'erreur individuelle.
Une référence intelligente pointe vers l'\cref{eq:frequence_cardiaque} située à la \cpageref{eq:frequence_cardiaque}.

\subsection{Structuration par paragraphes avec puces en ligne}
\label{subsec:paragraphes}

La classe personnalise la commande \texttt{\textbackslash paragraph} pour insérer une puce discrète et enchaîner le texte directement sur la même ligne :

\paragraph{Critère de complétude}
Le document doit couvrir l'ensemble des axes d'analyse de manière exhaustive et rigoureuse.

\paragraph{Critère d'accessibilité visuelle}
Les contrastes de couleurs, la taille des polices et la structure hiérarchique doivent garantir un confort de lecture optimal.

\subsection{Encadrés et boîtes visuelles (tcolorbox)}
\label{subsec:boites}

Le package \texttt{tcolorbox} est configuré pour enrichir la mise en page d'éléments d'accroche :

\begin{boiteSynthese}{Synthèse méthodologique}
  Ce modèle applique les principes de modularité, de compilation reproductible et de séparation stricte entre le contenu et la mise en forme.
  Toutes les fonctionnalités additionnelles (annotations rétractables, encadrés, géométries adaptatives) sont prêtes à l'emploi.
\end{boiteSynthese}

%------------------------------------------------------------------------------------
%  SECTION 2 : TABLEAUX PROFESSIONNELS
%------------------------------------------------------------------------------------
\section{Tableaux Comparatifs et Synthèse}
\label{sec:tableaux}

Pour obtenir un rendu typographique de qualité supérieure, les tableaux doivent impérativement respecter les règles du package \texttt{booktabs} :
jamais de barres verticales, une séparation claire par les règles horizontales (\texttt{\textbackslash toprule}, \texttt{\textbackslash midrule}, \texttt{\textbackslash bottomrule}), et un dimensionnement automatique via \texttt{tabularx}.

L'\cref{tab:hypothesesDiagnostiques} ainsi que l'\cref{tab:signesCliniques} illustrent la structure recommandée pour des tableaux médicaux ou scientifiques complexes.

% Inclusion propre via \input (et non \include qui casse la pagination d'une section)
%------------------------------------------------------------------------------------
%  TABLEAU DES HYPOTHÈSES DIAGNOSTIQUES (BOOKTABS & TABULARX PROFESSIONNEL)
%------------------------------------------------------------------------------------
\begin{table}[htbp]
  \centering
  \caption{Synthèse des hypothèses diagnostiques gynécologiques}
  \label{tab:hypothesesDiagnostiques}
  \small
  \begin{tabularx}{\textwidth}{>{\bfseries}p{3cm} X X}
    \toprule
    Diagnostic & Définition & Symptômes \\
    \midrule
    Kyste hémorragique du corps jaune &
    Kyste fonctionnel dépendant du cycle menstruel. &
    Dysménorrhée avec douleur unilatérale lors de la phase lutéale. \\
    \addlinespace
    Endométriome &
    Endomètre ectopique au niveau ovarien. &
    Dysménorrhée, douleurs pelviennes chroniques, dyspareunie profonde, troubles de la fertilité. \\
    \addlinespace
    Kyste dermoïde &
    Tératome, tumeur bénigne d'origine embryonnaire, fréquente entre $16$ et $55$ ans. &
    Asymptomatique ou légère douleur pelvienne, métrorragie, dyspareunie et compression éventuelle. \\
    \addlinespace
    Kyste para-ovarien &
    Kyste bénin avec pour origine le mésosalpinx. &
    Asymptomatique ou symptômes liés à une compression mécanique. \\
    \addlinespace
    Tumeur ovarienne &
    Cancer ovarien malin. Facteurs favorisants : $G_0 P_0$, ménarche précoce, cycles spontanés. &
    \emph{Stade précoce :} paucisymptomatique.\newline
    \emph{Stade avancé :} AEG, perte pondérale, asthénie. \\
    \bottomrule
  \end{tabularx}
\end{table}

\begin{table}[htbp]
  \centering
  \caption{Signes cliniques et remarques associées}
  \label{tab:signesCliniques}
  \small
  \begin{tabularx}{\textwidth}{>{\bfseries}p{3cm} X X}
    \toprule
    Diagnostic & Signes cliniques & Remarques \\
    \midrule
    Kyste hémorragique du corps jaune &
    \begin{itemize}[nosep, leftmargin=*]
      \item Palpation abdominale : sensible ;
      \item Toucher vaginal : masse annexielle palpée, mobilisation utérine douloureuse.
    \end{itemize} &
    Rupture du kyste $\rightarrow$ défense abdominale, signe du cri du Douglas. \\
    \addlinespace
    Endométriome &
    \begin{itemize}[nosep, leftmargin=*]
      \item Palpation abdominale : sensible ;
      \item Toucher vaginal : masse annexielle palpable, adhérences suspectées.
    \end{itemize} &
    Symptômes associés possibles (nycturie, troubles du transit). \\
    \addlinespace
    Kyste dermoïde &
    \begin{itemize}[nosep, leftmargin=*]
      \item Palpation abdominale : normale ;
      \item Toucher vaginal : masse annexielle potentiellement palpable.
    \end{itemize} &
    Torsion ovarienne = urgence chirurgicale $\rightarrow$ douleur aiguë et brutale. \\
    \addlinespace
    Kyste para-ovarien &
    \begin{itemize}[nosep, leftmargin=*]
      \item Palpation abdominale : normale ;
      \item Toucher vaginal : masse annexielle parfois palpable.
    \end{itemize} &
    Évolution le plus souvent asymptomatique. \\
    \addlinespace
    Tumeur ovarienne &
    \begin{itemize}[nosep, leftmargin=*]
      \item Palpation abdominale : signes de carcinose ;
      \item Toucher vaginal : masse annexielle dure, fixée.
    \end{itemize} &
    Diagnostic d'exclusion prioritaire en raison du pronostic vital. \\
    \bottomrule
  \end{tabularx}
\end{table}




%------------------------------------------------------------------------------------
%  SECTION 3 : GESTION DES FIGURES ET ILLUSTRATIONS
%------------------------------------------------------------------------------------
\section{Gestion des Figures et Illustrations}
\label{sec:illustrations}

Le modèle sépare les figures techniques conventionnelles des illustrations visuelles ou artistiques grâce à une liste dédiée.

\subsection{Figure standard avec sous-figures}
\label{subsec:figures_standards}

La \cref{fig:exemple_figures} présente une organisation de sous-figures alignées avec des légendes indépendantes.

\begin{figure}[htbp]
  \centering
  \begin{subfigure}[b]{0.45\textwidth}
    \centering
    \rule{0.8\textwidth}{3cm}
    \caption{Schéma technique A}
    \label{subfig:schema_a}
  \end{subfigure}
  \hfill
  \begin{subfigure}[b]{0.45\textwidth}
    \centering
    \rule{0.8\textwidth}{3cm}
    \caption{Schéma technique B}
    \label{subfig:schema_b}
  \end{subfigure}
  \caption{Exemple de figure double avec sous-légendes}
  \label{fig:exemple_figures}
\end{figure}

\subsection{Environnement Illustration personnalisé avec cadre décoratif}
\label{subsec:illustrations_dediees}

L'environnement \texttt{illustration} permet d'indexer les supports iconographiques dans une table distincte.
Le cadre personnalisé \texttt{cadreFigure} permet d'englober l'image d'un contour soigné aux coins arrondis :

\begin{illustration}[htbp]
  \centering
  \begin{cadreFigure}[width=0.7\textwidth]
    \rule{7cm}{3.5cm}
  \end{cadreFigure}
  \caption{Représentation anatomique ou graphique d'expérimentation}
  \label{ill:anatomie}
\end{illustration}

L'\cref{ill:anatomie} apparaît automatiquement dans la \emph{Liste des illustrations}.


% --- CHAPITRE II ---
\chapter{Méthodologie et Analyse Pratique}
\label{chap:methodologie}

%------------------------------------------------------------------------------------
%  SECTION 1 : ANNOTATIONS ET SUIVI DES MODIFICATIONS
%------------------------------------------------------------------------------------
\section{Annotations et Notes de Suivi}
\label{sec:todos}

Durant la phase de rédaction, il est essentiel de pouvoir annoter le document sans polluer le texte définitif.
L'option de classe \texttt{todoNote} active l'affichage des notes marginales et de la table des révisions.

\subsection{Notes dans la marge et surlignage}
\label{subsec:notes_marge}

La commande personnalisée \texttt{\textbackslash todoH} permet d'associer directement un surlignage coloré dans le corps du texte à une note en marge :

Dans cette phrase, nous identifions un point à relire avec attention \todoH{MarieStyle}{Emerald!30}{Vérifier les données statistiques}{les données cliniques présentées ici}.
De plus, une seconde remarque peut être formulée pour l'auteur \todoH{GuillaumeStyle}{Apricot!40}{Ajouter la référence manquante}{concernant l'algorithme d'imagerie médicale}.

\subsection{Gestion des figures provisoires}
\label{subsec:missing_figures}

Lorsqu'une figure n'est pas encore finalisée, la commande \texttt{\textbackslash missingfigure} réserve l'espace dans la mise en page sans bloquer la compilation :

\missingfigure[figwidth=\linewidth, figheight=4.5cm, figcolor=RoyalBlue!20]{Schéma de synthèse du protocole d'étude à insérer}


%------------------------------------------------------------------------------------
%  SECTION 2 : CITATIONS ET BIBLIOGRAPHIE
%------------------------------------------------------------------------------------
\section{Citations et Gestion Bibliographique}
\label{sec:bibliographie}

Le moteur de gestion bibliographique utilisé est \texttt{biblatex} associé au backend \texttt{biber}.
Le style sélectionné ici est Vancouver, adapté aux publications médicales et scientifiques.

\subsection{Exemples de citations en contexte}
\label{subsec:exemples_citations}

Différents modes de citation sont pris en charge selon la syntaxe de la phrase :
\begin{itemize}
  \item \textbf{Citation intégrée au texte (auteur-date ou sujet) :}
  Comme le démontrent \textcite{nieva2024impact}, l'évolution des recommandations influence directement les diagnostics précoces.
  \item \textbf{Citation globale ou parenthétique :}
  Plusieurs biomarqueurs sanguins prometteurs font l'objet d'évaluations multicentriques récentes~\parencite{burghaus2024}.
  \item \textbf{Citation avec numéro standard :}
  L'échographie transvaginale demeure un examen clé dans la détection des lésions ovariennes~\cite{kanti2024ultrasound}.
  \item \textbf{Sources web et ressources en ligne :}
  Les principes de communication interpersonnelle sont également documentés en ligne~\parencite{demetersante2022}.
\end{itemize}

Les liens hypertextes créés sont directement cliquables et renvoient vers la section bibliographique en fin de document.
De plus, la fonctionnalité \texttt{backref} affiche les pages précises où chaque document a été référencé.



%------------------------------------------------------------------------------------
%  ANNEXES (Numérotation alphabétique avec préfixe "Annexe")
%------------------------------------------------------------------------------------
\appendix
\chapter{Documents Complémentaires et Pièces Jointes}
\label{annexe:justificatifs}

Cette section démontre l'intégration de documents externes dans le rapport final.
La commande fournie par le paquet \texttt{pdfpages} permet d'insérer des comptes-rendus, attestations ou fiches administratives au format PDF tout en conservant le style de page :

\begin{singlespace}
\begin{verbatim}
% Exemple d'inclusion de document PDF multipage :
\includepdf[pages=-, pagecommand={\thispagestyle{plain}}]{documents/annexe.pdf}
\end{verbatim}
\end{singlespace}

%------------------------------------------------------------------------------------
%  POST-CONTENU (BACKMATTER : Listes de figures, tables et bibliographie)
%------------------------------------------------------------------------------------
\backmatter

\listoffigures
\listofillustrations
\listoftables

\printbibliography[heading=bibintoc, title={Références Bibliographiques}]

\end{document}
